\section{Introduction}

Financial markets, much like other sectors of the economy, have become more data-driven over recent decades.
Whereas investors historically relied on a narrow set of macroeconomic and accounting indicators to form return expectations, contemporary market participants have access to an unprecedented volume of information.
These include a vast array of alternative signals, including textual data, satellite imagery, and news narratives where the relationship with future asset returns is seldomly clear.
Nevertheless, survey evidence indicates that the demand for these unconventional data sources continues to grow rapidly \citep{lowenstein_sandler_2026}.
To compensate for the lack of guiding economic theory, investors increasingly deploy machine learning techniques to process this high-dimensional information \citep{cfa_invesco_2024}.
These statistical algorithms enable investors to discover empirical patterns organically, without imposing rigid prior restrictions on model specification or variable selection.
However, this dynamic process of broad information acquisition and data-driven belief updating is not adequately captured by existing asset-pricing literature.
Standard adaptive learning frameworks typically assume that agents form expectations by recursively updating a fixed, low-dimensional perceived law of motion.
While analytically convenient, such frameworks provide an inadequate description of modern investment behavior, wherein cognitively constrained agents must evaluate thousands of potential signals without committing to a definitive structural model from the outset.

We develop a model that embeds boundedly rational belief formation and statistical learning into an otherwise standard asset-pricing environment.
Agents in our model suspect that returns may be conditionally predictable from a high-dimensional set of observable predictors.
To exploit this perceived opportunity while managing information overload, they employ a learning rule based on the least absolute shrinkage and selection operator (LASSO) \citep{Tibshirani1996}.
Functionally, this mirrors a cognitive process of selective attention.
Agents first estimate a perceived law of motion (PLM) using least squares and then reduce the dimensionality of their forecasting rule by shrinking all coefficients below a certain threshold to zero, effectively ignoring weak signals.

Our main result is that the act of searching for predictability among many candidate signals can itself generate short-lived pockets of return predictability.
In equilibrium, these episodes are sparse and transient, in line with recent evidence on episodic short-horizon predictability \citep{Chinco2019, Farmer2023}.
The mechanism follows a simple life cycle driven by the interaction between belief updating and bounded memory.
Because agents heavily discount historical data and rely on recent observations to estimate their perceived law of motion, their learning heuristic occasionally selects predictors whose finite-sample correlations are temporarily large enough to survive the shrinkage threshold.
When investors update their beliefs and trade on this perceived predictability, their demand feeds back into prices, transforming the perceived signal into genuine but temporary return predictability.
Crucially, this predictability is transient precisely because of bounded memory: as new data arrive and the specific noisy observations that triggered the initial selection fade from the agents' effective memory, the estimated coefficients revert toward zero.
Consequently, the narrative loses its salience, and prices return to behavior consistent with unconditional unpredictability.

We reach this result by formalizing the mechanism in a Lucas-type exchange economy with a single risky asset \citep{Lucas1978}.
Under rational expectations, the price–dividend ratio is constant and returns are unpredictable.
We follow \citet{Adam2016} and depart from rational expectations by assuming that agents know the fundamental dividend and consumption process but not the process for returns.
Instead, they suspect that returns are conditionally predictable from a large set of observable signals.
Each period, investors estimate a PLM in which next period's return depends linearly on a set of predictors.
Because the number of candidate predictors exceeds the number of observation available, standard linear updating is infeasible, prompting the use of a soft-thresholding rule similar to the LASSO.
These beliefs then feed into equilibrium prices through the actual law of motion (ALM) governing returns.
In contrast to the linear PLM, the ALM is nonlinear and thus misspecified.
In this case, learning is unlikely to converge to the true law of motion.
Instead, the economy converges to a cross-moment consistent equilibrium (CMCE) in which agents' subjective forecasting moments coincide with those generated by the equilibrium itself.
Occasional shifts in attention to specific variables trigger predictable return episodes that persist only until new observations arrive and replace old observations in agents memory.

We test the model's predictions using a large sample of NYSE-listed stocks and a set of 180 news attention series constructed in \citet{Bybee2021}.
Our estimations confirm several key implications of the model.
First, we document that agents uncover episodic return predictability concentrated in a small and shifting set of signals, consistent with \citet{Chinco2019}.
Second, we show that the frequency of these episodes is related to the volatility of the return process.
Third, we find that our framework of selective attention and endogenous belief-price feedback can explain a significant fraction of the cross-sectional variation in stock returns.

The rest of the paper is structured as follows.
Section~\ref{sec:literature} discusses related literature.
Section~\ref{sec:model} presents the model environment.
Section~\ref{sec:rec_ls} analyzes the asset-pricing equilibrium under recursive least squares.
Section~\ref{sec:w_ls} analyzes the asset-pricing equilibrium under weighted least squares.
Section~\ref{sec:lasso} introduces LASSO learning and derives our main results.
Section~\ref{sec:empirical} estimates the model on a sample of daily stock returns.
Section~\ref{sec:conclusion} concludes.

